% Generated draft: requires mathematical and source review.
\documentclass[11pt]{article}
\usepackage[margin=25mm]{geometry}
\usepackage{fontspec}
\setmainfont{Latin Modern Roman}
\usepackage{amsmath,amssymb,mathtools,unicode-math}
\setmathfont{Latin Modern Math}
\usepackage{enumitem,needspace,xurl,hyperref}
\setlength{\parindent}{0pt}
\setlength{\parskip}{5pt}
\setlength{\emergencystretch}{3em}
\newcommand{\N}{\mathbb{N}}
\newcommand{\Z}{\mathbb{Z}}
\newcommand{\Q}{\mathbb{Q}}
\newcommand{\R}{\mathbb{R}}
\newcommand{\C}{\mathbb{C}}
\newcommand{\F}{\mathbb{F}}
\newcommand{\A}{\mathbb{A}}
\newcommand{\PP}{\mathbb{P}}
\newcommand{\E}{\mathbb{E}}
\newcommand{\eps}{\varepsilon}
\newcommand{\dd}{\,\mathrm{d}}
\begin{document}
\section*{UnsolvedMath ID: 700450}
\textbf{Source problem code:} catalogue entry
\subsection{Definitions and mathematical statement}
A group \(H\) satisfies the \emph{minimum condition on subgroups} if every descending chain of subgroups
\[
H_0\ge H_1\ge H_2\ge\cdots,\qquad H_i\le H,
\]
is eventually constant: there exists an integer \(n\ge0\) such that \(H_i=H_n\) for all \(i\ge n\). It satisfies the \emph{maximum condition on subgroups} if every ascending chain
\[
H_0\le H_1\le H_2\le\cdots,\qquad H_i\le H,
\]
is eventually constant in the same sense.

For subgroups \(A,B\le G\), write
\[
AB=\{ab:a\in A,\ b\in B\}.
\]
The hypothesis \(G=AB\) means that every element of \(G\) has such a representation. Neither normality of the factors nor uniqueness of the representation is assumed. No finiteness hypothesis is supplied.

\textbf{Problem.} Determine whether each of the following two assertions holds for every group \(G\) and all subgroups \(A,B\le G\) with \(G=AB\):
\begin{enumerate}
\item If both \(A\) and \(B\) satisfy the minimum condition on subgroups, then \(G\) satisfies the minimum condition on subgroups.
\item If both \(A\) and \(B\) satisfy the maximum condition on subgroups, then \(G\) satisfies the maximum condition on subgroups.
\end{enumerate}
The source asks these as separate minimum and maximum condition questions. Its statement begins with an attribution to B. Amberg, whereas its proposer metadata and background name N. S. Chernikov; their respective roles are unresolved in the supplied record.
\subsection{Short English statement}
If a group is the setwise product of two subgroups whose descending subgroup chains eventually stabilize, must its own descending subgroup chains stabilize? Ask the same question separately for ascending chains. The supplied attribution to B. Amberg conflicts with proposer metadata naming N. S. Chernikov.
\subsection{Sources}
\begin{itemize}
\item Supplied local UnsolvedMath registry record, permanent ID \(700450\), title \emph{Kourovka Notebook Problem 7.56}, exact source code \texttt{catalogue entry}; upstream revision \texttt{d371ee917ac9647fbb06cd6fc529deec6f3dda82}.
\item Supplied source attribution: \emph{The Kourovka Notebook (Unsolved Problems in Group Theory)}, 7th issue (1980), problem 7.56; also identified in the record as appearing in the 21st issue edition (2026), \texttt{arXiv:1401.0300v48}. The collection description names E. I. Khukhro and V. D. Mazurov as editors. These bibliographic details are supplied and have not been independently verified.
\item The statement explicitly names B. Amberg in parentheses. The background and proposer field instead name N. S. Chernikov, with proposed year 1980. The record does not explain the attribution difference.
\item The local registry labels the problem \texttt{open}. This is a supplied status label, not an independently verified assessment of either assertion. No separate literature assessment, verification date, or literature-source links are supplied.
\end{itemize}
\end{document}
